% theory_optimal_capital.tex
% Optimal financing channel selection, WACBA, and endogenous leverage dynamics

\section{Optimal Capital Structure and Financing Channel Selection}
\label{sec:optimal_capital}

Strategy employs five distinct financing channels to acquire BTC:
at-the-market (ATM) common equity issuance, convertible senior notes,
and three classes of perpetual preferred stock---convertible (STRK, 8\%),
non-convertible fixed-rate (STRF 10\%, STRD 10\%, STRE 10\% EUR-denominated),
and variable-rate (STRC, currently 11.5\%).
Each channel carries a different effective cost of BTC acquisition
that depends on the log-premium $\pi_t$ and the BTC spot price $S_t$.
This section derives the cost structure for each channel, establishes
the optimal financing rule, and introduces the Weighted Average Cost
of BTC Acquisition (WACBA).

\subsection{Cost of Each Financing Channel}

We define the \emph{effective cost of BTC acquisition} for channel $c$
as the rate $k_c(\pi_t, S_t)$ such that an incremental dollar raised
through channel $c$ adds BTC at an all-in cost of $k_c$ per unit,
expressed as an annualized percentage drag on BTC-per-share accretion.

\subsubsection{ATM Common Equity}

When the firm issues $\Delta N$ new shares at market price $P_t$,
it raises $P_t \Delta N$ dollars and purchases
$\Delta H = P_t \Delta N / S_t$ BTC.
The net effect on BTC per share is:
\begin{equation}
  \Delta B = \frac{H_t + \Delta H}{N_t^{\mathrm{sh}} + \Delta N} - B_t
  = \frac{B_t(e^{\pi_t} \cdot \text{ILE}_t^{-1} - 1)}{1 + \Delta N / N_t^{\mathrm{sh}}}\cdot\frac{\Delta N}{N_t^{\mathrm{sh}}},
  \label{eq:atm_dilution}
\end{equation}
where the simplification uses
$P_t / S_t = e^{\pi_t}(H_t S_t - D_t) / (N_t^{\mathrm{sh}} S_t)
= e^{\pi_t} B_t (1 - D_t / (H_t S_t))$.

\begin{proposition}[ATM Accretion Threshold]
\label{prop:atm_threshold}
ATM equity issuance is BTC-per-share accretive if and only if
\begin{equation}
  \pi_t > \log\!\bigl(\mathrm{ILE}_t\bigr) = \log\!\left(\frac{H_t S_t}{H_t S_t - D_t}\right).
  \label{eq:atm_accretive}
\end{equation}
When accretive, the marginal cost of BTC acquisition via ATM is
\begin{equation}
  k_{\mathrm{ATM}}(\pi_t) = -\mu_{H/N}^{\mathrm{marginal}}
  = -\frac{e^{\pi_t}/\mathrm{ILE}_t - 1}{e^{\pi_t}/\mathrm{ILE}_t} < 0,
  \label{eq:atm_cost}
\end{equation}
i.e., negative cost (value-creating). When $\pi_t \le \log(\mathrm{ILE}_t)$,
ATM issuance is dilutive and the firm optimally avoids this channel.
\end{proposition}

\begin{proof}
Write $P_t = e^{\pi_t}(A_t - D_t)/N_t^{\mathrm{sh}}$ and
$B_t = H_t / N_t^{\mathrm{sh}}$.
Each new share issued at price $P_t$ purchases $P_t / S_t$ BTC.
The issuance is accretive iff the BTC purchased per new share exceeds
the existing BTC per share:
\[
  \frac{P_t}{S_t} > B_t
  \iff
  \frac{e^{\pi_t}(H_t S_t - D_t)}{N_t^{\mathrm{sh}} S_t} > \frac{H_t}{N_t^{\mathrm{sh}}}
  \iff
  e^{\pi_t} > \frac{H_t S_t}{H_t S_t - D_t} = \mathrm{ILE}_t.
\]
The cost expression follows by computing the fractional BTC-per-share
gain per dollar deployed.
\end{proof}

\begin{remark}
With current parameters (Table~\ref{tab:calibrated}), $\mathrm{ILE}_t = 1.488$
so the accretion threshold is $\pi_t > \log(1.488) \approx 0.397$.
The current premium $\pi_0 = 0.242$ is \emph{below} this threshold,
indicating that ATM equity issuance would currently be mildly dilutive,
consistent with the observed shift toward preferred stock issuance.
\end{remark}

\subsubsection{Convertible Senior Notes}

A convertible note with coupon $c_j$, conversion price $K_j$, and maturity $T_j$
has effective cost comprising two components:
\begin{enumerate}
  \item \textbf{Coupon cost:} The ongoing interest payment $c_j$ per annum
    on principal $D_j$.
  \item \textbf{Expected dilution at conversion:} If $P_{T_j} > K_j$,
    the note converts into $D_j / K_j$ shares, diluting existing shareholders.
\end{enumerate}

The all-in annualized cost is:
\begin{equation}
  k_{\mathrm{conv},j} = c_j + \frac{1}{T_j - t}\,
  \mathbb{E}^{\mathbb{P}}\!\left[
    \left(\frac{D_j}{K_j}\cdot\frac{B_{T_j}}{1 + D_j/(K_j N_{T_j}^{\mathrm{sh}})}
    - \frac{D_j}{S_{T_j}}\right)^{\!+}
    \cdot \frac{S_{T_j}}{D_j}
  \right],
  \label{eq:conv_cost}
\end{equation}
where the expectation captures the path-dependent dilution loss relative
to outright BTC purchase.

\begin{remark}
Strategy's convertible notes carry exceptionally low coupons
(weighted average 0.42\%), making them the cheapest debt channel.
The 0\% notes due 2029 (\$3B, conversion at \$672.40) and 2030
(\$2B, conversion at \$433.43) dominate the outstanding convertible portfolio.
At current MSTR price levels, notes with conversion prices below \$250
are deep in-the-money and should be modeled as quasi-equity.
\end{remark}

\subsubsection{Non-Convertible Perpetual Preferred Stock}

For non-convertible preferred (STRF, STRD, STRE), the cost is
straightforward:
\begin{equation}
  k_{\mathrm{pref},j} = \frac{d_j \cdot L_j}{\text{Proceeds}_j},
  \label{eq:pref_cost}
\end{equation}
where $d_j$ is the dividend rate, $L_j$ the liquidation preference,
and $\text{Proceeds}_j$ the net issuance proceeds per share.

\begin{center}
\begin{tabular}{lcccc}
\toprule
Ticker & Div.\ Rate & IPO Price & Stated Amount & Effective Cost \\
\midrule
STRF & 10.0\% & \$85.00 & \$100.00 & 11.76\% \\
STRD & 10.0\% & \$85.00 & \$100.00 & 11.76\% \\
STRE & 10.0\% & \EUR{80.00} & \EUR{100.00} & 12.50\% \\
STRC & 9.6\%$^*$ & \$90.00 & \$100.00 & 10.67\%--12.78\%$^\dagger$ \\
\bottomrule
\end{tabular}
\end{center}
{\small $^*$Initial rate; current rate 11.5\%.
$^\dagger$Range reflects variable rate adjustment.}

\medskip
These instruments carry no dilution risk to common shareholders
but impose a perpetual fixed obligation.
The annual dividend burden is approximately \$976.6M.

\subsubsection{Convertible Perpetual Preferred (STRK)}

STRK (8\% cumulative, convertible at 0.1 MSTR shares per STRK share,
effective conversion price \$1{,}000) combines preferred dividend cost
with conversion optionality:
\begin{equation}
  k_{\mathrm{STRK}} = d_{\mathrm{STRK}} + \frac{1}{T_{\mathrm{eff}}}
  \mathbb{E}^{\mathbb{P}}\!\left[
    \mathbf{1}_{\{P_T > K_{\mathrm{STRK}}\}}
    \cdot \text{dilution cost}
  \right],
  \label{eq:strk_cost}
\end{equation}
where $d_{\mathrm{STRK}} = 8\%/\$80 = 10\%$ effective dividend yield
and $K_{\mathrm{STRK}} = \$1{,}000$ per MSTR share.
The conversion component is currently deep out-of-the-money,
making STRK's effective cost approximately 10\%, comparable to
non-convertible preferred but with a \$21B ATM capacity that
provides substantial future issuance optionality.

\subsection{Optimal Financing Rule}

\begin{proposition}[Optimal Channel Selection]
\label{prop:optimal_channel}
Let $\pi^*_{\mathrm{ATM}} = \log(\mathrm{ILE}_t)$ denote the ATM accretion
threshold and $\pi^*_{\mathrm{conv}}$ the premium level at which
convertible issuance becomes cheaper than preferred.
The cost-minimizing financing rule is:
\begin{equation}
  c^*(\pi_t) =
  \begin{cases}
    \text{ATM equity} & \text{if } \pi_t > \pi^*_{\mathrm{ATM}},
      \quad k_{\mathrm{ATM}} < 0, \\[4pt]
    \text{Convertible bonds} & \text{if } \pi^*_{\mathrm{conv}} < \pi_t
      \le \pi^*_{\mathrm{ATM}},
      \quad k_{\mathrm{conv}} \approx 0.4\text{--}2.3\%, \\[4pt]
    \text{Preferred stock} & \text{if } \pi_t \le \pi^*_{\mathrm{conv}},
      \quad k_{\mathrm{pref}} \approx 10\text{--}12.5\%.
  \end{cases}
  \label{eq:optimal_channel}
\end{equation}
\end{proposition}

\begin{proof}
The ordering follows directly from the cost functions derived above.
When $\pi_t > \pi^*_{\mathrm{ATM}}$, ATM issuance has negative cost
(accretive) and strictly dominates all other channels.
Below this threshold, convertible bonds offer the lowest positive cost
due to their near-zero coupons, provided the premium is sufficient
to attract investor demand for equity upside.
When the premium is low or negative, only the unconditional claim
of preferred dividends can attract capital, at the highest cost.

The thresholds are:
\begin{align}
  \pi^*_{\mathrm{ATM}} &= \log(\mathrm{ILE}_t)
  = \log\!\left(\frac{A_t}{A_t - D_t}\right), \\
  \pi^*_{\mathrm{conv}} &\approx \log(\mathrm{ILE}_t)
  - \frac{k_{\mathrm{pref}} - \bar{c}_{\mathrm{conv}}}
         {\gamma_{\mathrm{demand}}},
\end{align}
where $\gamma_{\mathrm{demand}} > 0$ parameterizes investor demand
sensitivity to premium for convertible issuance.
\end{proof}

\begin{remark}[Empirical Validation]
Strategy's observed behavior closely matches the optimal rule:
\begin{itemize}
  \item In Q4 2024 ($\pi_t \approx 1.5$--$2.0$, well above threshold):
    massive ATM equity issuance (\$15B+) and convertible issuance (\$5B).
  \item In Q1--Q2 2025 ($\pi_t$ declining toward $0.5$--$1.0$):
    shift to STRK and STRF preferred issuance.
  \item In H2 2025--Q1 2026 ($\pi_t \approx 0.2$--$0.5$, near/below threshold):
    heavy STRC, STRD issuance; STRE launched for EUR diversification.
    ATM equity issuance slowed significantly (\$8.1B remaining capacity).
\end{itemize}
The progression from equity $\to$ convertibles $\to$ preferred
as $\pi_t$ declines is consistent with Proposition~\ref{prop:optimal_channel}.
\end{remark}

\subsection{Weighted Average Cost of BTC Acquisition (WACBA)}

\begin{definition}[WACBA]
\label{def:wacba}
The Weighted Average Cost of BTC Acquisition is
\begin{equation}
  \mathrm{WACBA}_t = \sum_{c \in \mathcal{C}} w_c(t)\, k_c(\pi_t, S_t),
  \label{eq:wacba}
\end{equation}
where $\mathcal{C} = \{\mathrm{ATM}, \mathrm{conv}_1, \ldots,
\mathrm{conv}_J, \mathrm{STRK}, \mathrm{STRF}, \mathrm{STRC},
\mathrm{STRD}, \mathrm{STRE}\}$ is the set of financing channels,
$w_c(t) = \Delta H_c(t) / \sum_{c'} \Delta H_{c'}(t)$ is the share
of BTC acquired through channel $c$ over a trailing window, and
$k_c(\pi_t, S_t)$ is the effective cost defined in
\S\ref{sec:optimal_capital}.1.
\end{definition}

WACBA is analogous to the Weighted Average Cost of Capital (WACC)
in corporate finance, but measures the blended cost of the firm's
BTC acquisition program rather than its capital structure per se.

\begin{proposition}[WACBA Minimization]
\label{prop:wacba_min}
The management optimization problem is:
\begin{equation}
  \min_{\{w_c\}_{c \in \mathcal{C}}} \mathrm{WACBA}_t
  \quad \text{subject to} \quad
  \begin{cases}
    w_c \ge 0, \quad \sum_c w_c = 1, \\
    w_c \le \bar{w}_c(t) \quad \text{(capacity constraints)}, \\
    \sum_c D_c(t) \le \bar{D}(t) \quad \text{(leverage constraint)}, \\
    \sum_c d_c w_c \le \bar{d}(t) \quad \text{(dividend coverage)}.
  \end{cases}
  \label{eq:wacba_min}
\end{equation}
\end{proposition}

The capacity constraints $\bar{w}_c(t)$ reflect remaining ATM program
authorizations: \$8.1B common equity, \$20.3B STRK, \$3.6B STRC,
\$4.2B STRD, \$2.1B STRF.
The leverage constraint ensures $\mathrm{ILE}_t$ does not exceed
a prudential bound.
The dividend coverage constraint ensures the annual preferred dividend
obligation ($\approx$\$976.6M) remains serviceable from BTC holdings
or cash reserves (\$2.25B, providing 2.5 years of coverage).

\subsection{Endogenous Leverage Dynamics}
\label{subsec:endo_leverage}

The interaction between premium, financing, and leverage creates
procyclical dynamics.

\begin{proposition}[Procyclical Leverage]
\label{prop:procyclical}
Under the structural model of Section~2, leverage dynamics are procyclical:
\begin{enumerate}
  \item \textbf{Bull market} ($S_t$ rising, $\pi_t$ high):
    \begin{itemize}
      \item ATM equity issuance is accretive ($\pi_t > \log\mathrm{ILE}_t$).
      \item Equity issuance increases $N_t^{\mathrm{sh}}$ and raises cash
        for BTC purchases, increasing $H_t$.
      \item Rising $S_t$ increases $A_t = H_t S_t$ faster than $D_t$,
        so $\mathrm{ILE}_t = A_t/(A_t - D_t) \to 1$.
      \item Leverage \emph{falls}.
    \end{itemize}

  \item \textbf{Bear market} ($S_t$ falling, $\pi_t$ low):
    \begin{itemize}
      \item ATM equity issuance becomes dilutive.
      \item Financing shifts to preferred stock (no common dilution
        but perpetual obligations).
      \item Falling $S_t$ decreases $A_t$ while $D_t$ is fixed,
        so $\mathrm{ILE}_t \to \infty$ as $A_t \to D_t$.
      \item Leverage \emph{rises}.
    \end{itemize}
\end{enumerate}
\end{proposition}

\begin{proof}
Differentiating $\mathrm{ILE}_t = A_t/(A_t - D_t)$ with respect to $S_t$:
\[
  \frac{\partial\,\mathrm{ILE}_t}{\partial S_t}
  = \frac{-D_t H_t}{(H_t S_t - D_t)^2} < 0.
\]
So ILE is decreasing in $S_t$: leverage falls as BTC price rises
and rises as BTC price falls. The financing channel switch
(Proposition~\ref{prop:optimal_channel}) amplifies this:
in bull markets the firm issues equity (deleveraging),
while in bear markets it issues debt-like preferred (adding obligations).
\end{proof}

This procyclicality connects to the margin spiral mechanism of
\citet{brunnermeier2009market}. In their framework, declining asset
values tighten funding constraints, forcing deleveraging that further
depresses prices. Strategy's variant is distinct: it does not face
margin calls on BTC (which is unlevered on-balance-sheet), but the
\emph{inability to issue accretive equity} in bear markets forces
reliance on expensive preferred financing, degrading the capital
structure and increasing the BTC price level required for future
accretive equity issuance.

\begin{corollary}[Leverage Spiral Condition]
\label{cor:leverage_spiral}
A destabilizing leverage spiral occurs if the preferred dividend burden
grows faster than BTC asset income. Formally, instability arises when
\begin{equation}
  \frac{d}{dt}\!\left(\sum_j d_j L_j\right) >
  \frac{d}{dt}\!\left(\mu_S \cdot H_t S_t\right),
  \label{eq:leverage_spiral}
\end{equation}
where $\mu_S$ is the expected BTC return. At current parameters,
with \$976.6M annual preferred dividends, expected BTC return
$\mu_S = 11.1\%$, and BTC assets of \$56.25B, the right-hand side
(\$6.24B) far exceeds the left-hand side, providing substantial buffer.
However, a sustained BTC decline to $S_t \approx \$15{,}000$--$\$20{,}000$
would close this gap.
\end{corollary}